\section*{Chapitre \ref{ch:b1:continuity} --- Limites et continuité}

\begin{solution}{exo:b1:continuity:1}
Prenons $u_n = \frac{1}{2\pi n + \pi/2}$ et $v_n = \frac{1}{2\pi n}$~:
toutes deux tendent vers $0^+$, et pourtant $\sin\frac{1}{u_n} = 1$ et
$\sin\frac{1}{v_n} = 0$. Deux suites, deux limites différentes des
images~: d'après le \cref{thm:b1:continuity:seqchar}, pas de limite en
$0^+$.

$x \sin\frac1x$~: encadrée par $\abs{x\sin\frac1x} \leq \abs x \to 0$,
donc la limite en $0$ existe et vaut $0$.
\end{solution}

\begin{solution}{exo:b1:continuity:2}
$f = \lfloor\cdot\rfloor$ est \omterm{def:b1:continuity:continuous}{continue} sur $\R \setminus \Z$
(localement constante) et discontinue en chaque $n \in \Z$~: limite à
gauche $n - 1$, valeur $n$.

$g(x) = x - \lfloor x\rfloor$~: mêmes points de discontinuité
(l'identité est \omterm{def:b1:continuity:continuous}{continue}, donc $g$ hérite des sauts de $f$)~; en $n \in
\Z$, la limite à gauche vaut $1 \neq 0 = g(n)$.

$h(x) = \lfloor x\rfloor + (x - \lfloor x\rfloor)^2$~: sur $\intco{n}{n
+ 1}$, $h(x) = n + (x - n)^2$, \omterm{def:b1:continuity:continuous}{continue}~; en $x = n$, la limite à
gauche vaut $(n-1) + 1 = n = h(n)$~: les sauts se compensent. $h$ est
\omterm{def:b1:continuity:continuous}{continue} sur $\R$ (et strictement croissante).
\end{solution}

\begin{solution}{exo:b1:continuity:3}
$P(x) = x^5 - 3x + 1$~: $P(-2) = -32 + 6 + 1 = -25 < 0$~; $P(0) = 1 >
0$~; $P(1) = -1 < 0$~; $P(2) = 27 > 0$. Trois changements de signe sur
les segments disjoints $\intcc{-2}{0}$, $\intcc{0}{1}$, $\intcc{1}{2}$~:
d'après le \cref{thm:b1:continuity:ivt}, au moins trois racines. (Étant
de degré $5$, $P$ en a au plus cinq~; une étude des variations
montrerait qu'il y en a exactement trois.)
\end{solution}

\begin{solution}{exo:b1:continuity:4}
Soit $P = a_{2m+1}X^{2m+1} + \dots$ avec $a_{2m+1} > 0$ (sinon
remplacer $P$ par $-P$). En factorisant le terme dominant, $P(x) =
a_{2m+1} x^{2m+1}\bigl(1 + o(1)\bigr)$ quand $x \to \pm\infty$~: donc
$P(x) \to +\infty$ en $+\infty$ et $-\infty$ en $-\infty$. Choisissons
$a$ tel que $P(a) < 0$ et $b$ tel que $P(b) > 0$~: le \omterm{thm:b1:continuity:ivt}{théorème des
valeurs intermédiaires} sur $\intcc{a}{b}$ fournit une racine.
\end{solution}

\begin{solution}{exo:b1:continuity:5}
Soit $g(x) = f(x) - x$, \omterm{def:b1:continuity:continuous}{continue} sur $\intcc{0}{1}$. Comme $f$ est à
valeurs dans $\intcc{0}{1}$~: $g(0) = f(0) \geq 0$ et $g(1) = f(1) - 1
\leq 0$. D'après le \cref{thm:b1:continuity:ivt}, $g(c) = 0$ pour un
certain $c$~: un point fixe.

Nécessité des hypothèses~: sur $\intcc{0}{1}$, l'\omterm{def:b1:logic:map}{application}
discontinue $f(x) = 1$ pour $x \leq \frac12$, $f(x) = 0$ pour $x >
\frac12$ n'a pas de point fixe~; sur l'\omterm{prop:b1:reals:intervals}{intervalle} $\intoo{0}{1}$ (qui
n'est pas un segment), $f(x) = \frac x2$ est \omterm{def:b1:continuity:continuous}{continue} à valeurs dans
$\intoo{0}{1}$ sans point fixe (le candidat $0$ manque)~; sur $\R$,
$f(x) = x + 1$.
\end{solution}

\begin{solution}{exo:b1:continuity:6}
$g(x) = \cos x - x$ est \omterm{def:b1:continuity:continuous}{continue}, $g(0) = 1 > 0$, $g(1) = \cos 1 -
1 < 0$~: une solution existe dans $\intoo{0}{1}$
(\cref{thm:b1:continuity:ivt}). Unicité~: $g$ est strictement
décroissante sur $\R$ --- pour $x \leq 0$, $g(x) \geq 1 - x > 0$ ne
s'annule de toute façon pas~; et $g'(x) = -\sin x - 1 \leq 0$ avec
égalité seulement en des points isolés ($x \equiv -\frac\pi2 \bmod
2\pi$), donc $g$ est strictement décroissante (\cref{ch:b1:derivative}~;
autrement~: sur $\intcc{0}{1}$, $\cos$ est strictement décroissante et
$-x$ aussi, donc $g$ l'est). Une fonction strictement monotone s'annule
au plus une fois.
\end{solution}

\begin{solution}{exo:b1:continuity:7}
Fixons $M = f(0) + 1$. Il existe $A > 0$ tel que $f(x) \geq M$ pour
$\abs x \geq A$ (définition des deux limites infinies~; prendre le plus
grand des deux seuils). Sur le segment $\intcc{-A}{A}$, le \omterm{thm:b1:continuity:evt}{théorème des
bornes atteintes} (\cref{thm:b1:continuity:evt}) fournit $c$ tel que
$f(c) = \inf_{\intcc{-A}{A}} f \leq f(0)$. Pour $\abs x \geq A$~: $f(x)
\geq M > f(0) \geq f(c)$. Donc $f(c)$ est le minimum global.
\end{solution}

\begin{solution}{exo:b1:continuity:8}
Sur le segment $\intcc{0}{T}$, $f$ est bornée et atteint ses bornes
(\cref{thm:b1:continuity:evt})~; par périodicité, ce sont les bornes
sur $\R$ tout entier, toujours atteintes.

Soit $g(x) = f(x + \frac T2) - f(x)$, \omterm{def:b1:continuity:continuous}{continue}. Alors
\[
g(0) + g\bigl(\tfrac T2\bigr)
= \bigl(f(\tfrac T2) - f(0)\bigr) + \bigl(f(T) - f(\tfrac T2)\bigr)
= f(T) - f(0) = 0 :
\]
$g(0)$ et $g(\frac T2)$ sont de signes opposés (ou bien l'un des deux
s'annule), donc le \omterm{thm:b1:continuity:ivt}{théorème des valeurs intermédiaires} sur
$\intcc{0}{T/2}$ donne $c$ tel que $g(c) = 0$, c'est-à-dire $f(c +
\frac T2) = f(c)$.
\end{solution}

\begin{solution}{exo:b1:continuity:9}
D'abord l'inégalité~: pour $0 \leq y \leq x$,
\[
\bigl(\sqrt y + \sqrt{x - y}\bigr)^2 = x + 2\sqrt{y(x-y)} \geq x,
\]
donc $\sqrt x \leq \sqrt y + \sqrt{x - y}$, c'est-à-dire $\sqrt x -
\sqrt y \leq \sqrt{x - y}$. D'où $\abs{\sqrt x - \sqrt y} \leq
\sqrt{\abs{x - y}}$ pour tous $x, y \geq 0$.

\omterm{def:b1:continuity:uniform}{Continuité uniforme}~: soit $\varepsilon > 0$~; prenons $\delta =
\varepsilon^2$~; alors $\abs{x - y} \leq \delta$ entraîne $\abs{\sqrt x
- \sqrt y} \leq \sqrt\delta = \varepsilon$.

Pas lipschitzienne près de $0$~: $\frac{\sqrt x - \sqrt 0}{x - 0} =
\frac{1}{\sqrt x} \to +\infty$ quand $x \to 0^+$, donc aucune constante
$k$ ne peut dominer tous les taux d'accroissement.
\end{solution}

\begin{solution}{exo:b1:continuity:10}
Supposons $f$ \omterm{def:b1:logic:inj}{injective}, \omterm{def:b1:continuity:continuous}{continue}, et non strictement monotone. Alors
il existe $a < b < c$ dans $I$ tels que $f(b)$ ne soit pas compris
entre $f(a)$ et $f(c)$ --- en effet, si pour \emph{tous} les triplets
la valeur médiane était comprise entre les deux extrêmes, $f$ serait
monotone (comparer deux paires quelconques~; brève vérification par
cas). Disons $f(b) > \max(f(a), f(c))$ (l'autre cas est symétrique~:
remplacer $f$ par $-f$). Choisissons $v$ tel que $\max(f(a), f(c)) < v
< f(b)$. Par le \omterm{thm:b1:continuity:ivt}{théorème des valeurs intermédiaires} appliqué sur
$\intcc{a}{b}$ puis sur $\intcc{b}{c}$, il existe $u_1 \in
\intoo{a}{b}$ et $u_2 \in \intoo{b}{c}$ tels que $f(u_1) = v = f(u_2)$~:
deux points distincts de même image, ce qui contredit l'\omterm{def:b1:logic:inj}{injectivité}.
\end{solution}

\begin{solution}{exo:b1:continuity:11}
$f(0) = f(0+0) = 2f(0)$ donne $f(0) = 0$~; $f(-x) = -f(x)$ vient de $0 =
f(x - x)$. Posons $\alpha = f(1)$. Par récurrence~: $f(n) = n\alpha$
pour $n \in \N$, puis pour $n \in \Z$ par imparité. Pour $q \in \N^*$~:
$q\,f(\frac pq) = f(p) = p\alpha$ (ajouter $\frac pq$ à lui-même $q$
fois), donc $f(\frac pq) = \alpha\frac pq$~: $f = \alpha\,
\mathrm{id}$ sur $\Q$.

Soient maintenant $x \in \R$ et $(r_n)$ une suite de rationnels telle
que $r_n \to x$ (\omterm{def:b1:topology:dense}{densité}, \cref{thm:b1:reals:density}, appliquée dans
des \omterm{prop:b1:reals:intervals}{intervalles} emboîtés~; ou bien $r_n = \frac{\lfloor
nx\rfloor}{n}$). Par \omterm{def:b1:continuity:continuous}{continuité}~: $f(x) = \lim f(r_n) = \lim \alpha r_n
= \alpha x$.
\end{solution}

\begin{solution}{exo:b1:continuity:12}
Soit $\varepsilon > 0$. Par la limite en $+\infty$, il existe $A$ tel
que $\abs{f(x) - \ell} \leq \frac\varepsilon2$ pour $x \geq A$~; d'où,
pour $x, y \geq A$~: $\abs{f(x) - f(y)} \leq \varepsilon$ (aucune
proximité n'est requise).

Sur le segment $\intcc{0}{A + 1}$, le \omterm{thm:b1:continuity:heine}{théorème de Heine}
(\cref{thm:b1:continuity:heine}) donne $\delta_0 > 0$ pour ce
$\varepsilon$~; posons $\delta = \min(\delta_0, 1)$.

Prenons maintenant $x, y \geq 0$ quelconques avec $\abs{x - y} \leq
\delta$, disons $x \leq y$. Si $y \leq A + 1$~: tous deux sont dans le
segment, et le $\delta_0$ de Heine s'applique. Sinon $y > A + 1$, et
alors $x \geq y - 1 > A$~: tous deux sont dans $\intco{A}{+\infty}$, où
l'argument de la limite s'applique. Dans les deux cas $\abs{f(x) -
f(y)} \leq \varepsilon$~: \omterm{def:b1:continuity:uniform}{continuité uniforme}.
\end{solution}

\begin{solution}{pb:b1:continuity:1}
\textbf{1.} Pour $n \in \N$~: $f(nx) = n f(x)$ par récurrence
($f((n+1)x) = f(nx) + f(x)$). De plus $f(0) = 2f(0)$ donne $f(0) =
0$, et $0 = f(x - x) = f(x) + f(-x)$ donne l'imparité, donc $f(nx) =
nf(x)$ pour $n \in \Z$. Pour $r = \frac pq$~: $q\,f\bigl(\tfrac
pq x\bigr) = f(px) = p\,f(x)$, donc $f(rx) = r f(x)$~: $f$ est
$\Q$-linéaire.

\textbf{2.} L'additivité donne, pour tous $x$ et $h$~:
\[
f(x + h) - f(x) = f(h) = f(x_0 + h) - f(x_0) .
\]
Quand $h \to 0$, le membre de droite tend vers $0$ par \omterm{def:b1:continuity:continuous}{continuité} en
$x_0$~; donc $f(x + h) \to f(x)$~: \omterm{def:b1:continuity:continuous}{continuité} en tout $x$. Alors
l'\cref{exo:b1:continuity:11} donne $f(x) = f(1)\,x$.

\textbf{3.} Deux fonctions additives \omterm{def:b1:continuity:continuous}{continues} sont de la forme
$cx$ et $c'x$ (question 2)~; si elles coïncident sur le \omterm{def:b1:structures:subgroup}{sous-groupe}
\omterm{def:b1:topology:dense}{dense} $\Z + \sqrt2\,\Z$ (\cref{exo:b1:reals:9}), alors $c\,g = c'g$
pour un certain $g \neq 0$ de ce \omterm{def:b1:structures:subgroup}{sous-groupe}~: $c = c'$, les fonctions
sont égales. Sans \omterm{def:b1:continuity:continuous}{continuité}~: la $\Q$-linéarité ne lie que les valeurs
aux $\Q$-combinaisons, et $1, \sqrt2$ sont $\Q$-indépendants ($\sqrt2
\notin \Q$), donc $f(1) = 0$, $f(\sqrt2) = 1$ est cohérent et impose,
sur le \omterm{def:b1:structures:subgroup}{sous-groupe},
\[
f(m + n\sqrt2) = m\,f(1) + n\,f(\sqrt2) = n .
\]

\textbf{4.} Pour $t \in \intcc{0}{\ell}$~: $a + t \in
\intcc{a}{b}$, donc $f(t) = f(a + t) - f(a) \leq M - f(a) =: M'$.

\textbf{5.} Pour $t \in \intcc{0}{\ell}$, on a aussi $\ell - t \in
\intcc{0}{\ell}$, et l'additivité donne $f(t) = f(\ell) - f(\ell
- t) \geq f(\ell) - M'$. D'où $\abs{f} \leq C$ sur
$\intcc{0}{\ell}$ avec $C = \max\bigl(\abs{M'}, \abs{f(\ell) -
M'}\bigr)$.

\textbf{6.} Pour $t \in \intcc{0}{\ell/n}$~: $nt \in
\intcc{0}{\ell}$ et $f(t) = \frac{f(nt)}{n}$ (question 1), donc
$\abs{f(t)} \leq \frac Cn$. Soit $\varepsilon > 0$~; choisissons $n >
\frac C\varepsilon$~: pour $0 \leq h \leq \frac\ell n$,
$\abs{f(h)} \leq \varepsilon$, et pour $h$ négatif on utilise $f(h) =
-f(-h)$. Ainsi $f(h) \to 0 = f(0)$ quand $h \to 0$~: \omterm{def:b1:continuity:continuous}{continuité} en $0$,
donc partout (question 2), donc $f(x) = f(1)x$.

\textbf{7.} Si $f$ est croissante sur $\intcc{a}{b}$, alors
$f(a) \leq f(x) \leq f(b)$ sur cet \omterm{prop:b1:reals:intervals}{intervalle}~: bornée, donc linéaire
par la question 6, $f(x) = cx$~; et $c(b - a) = f(b) - f(a) \geq 0$
impose $c \geq 0$.

\textbf{8.} (a)$\Rightarrow$(b)$\Rightarrow$(c)~: trivial.
(c)$\Rightarrow$(a)~: question 2. (a)$\Rightarrow$(d)~: une fonction
linéaire est monotone partout. (d)$\Rightarrow$(e)~: une fonction
monotone sur $\intcc{a}{b}$ y est bornée par ses valeurs aux
extrémités. (e)$\Rightarrow$(a)~: questions 4 à 6. Les cinq
\omterm{def:b1:logic:statement}{assertions} sont équivalentes --- c'est l'échelle de régularité.

\textbf{9.} La question 4 n'a utilisé que la majoration $M$~; la
question 5 a \emph{déduit} la minoration de la majoration via la
réflexion $f(t) = f(\ell) - f(\ell - t)$~; la question 6 a ensuite
tourné sur $\abs f \leq C$. Donc~: additive et \emph{majorée} sur un
\omterm{prop:b1:reals:intervals}{intervalle} non dégénéré entraîne déjà linéaire. Pour une minoration,
appliquer ceci à $-f$ (additive, majorée). Barreau (e) le plus fort~:
une borne unilatérale sur un seul \omterm{prop:b1:reals:intervals}{intervalle} suffit.

\textbf{10.} Si $\frac{f(x)}{x}$ valait une même constante $c$ pour
tout $x \neq 0$, $f$ serait linéaire~; il existe donc $u, v$ non nuls
tels que $\frac{f(u)}u \neq \frac{f(v)}v$, c'est-à-dire $u f(v) - v
f(u) \neq 0$~: le déterminant des vecteurs $(u, f(u))$, $(v,
f(v))$ est non nul, et ils engendrent $\R^2$.

\textbf{11.} Pour $r, s \in \Q$~: $f(ru + sv) = r f(u) + s f(v)$
(question 1 deux fois, plus l'additivité), donc le graphe contient
\[
\bigl(ru + sv,\; r f(u) + s f(v)\bigr)
= r\,(u, f(u)) + s\,(v, f(v)) .
\]
Étant donnés $(x_0, y_0)$ et $\varepsilon > 0$~: le système $2 \times
2$ $a(u, f(u)) + b(v, f(v)) = (x_0, y_0)$ a une (unique) solution
réelle $(a, b)$ puisque le déterminant est non nul. Choisissons des
rationnels $r_n \to a$, $s_n \to b$~: alors $r_n(u, f(u)) +
s_n(v, f(v)) \to (x_0, y_0)$ coordonnée par coordonnée, et chacun de
ces points est sur le graphe~: le graphe est \omterm{def:b1:topology:dense}{dense} dans $\R^2$.

\textbf{12.} Soit $I$ un \omterm{prop:b1:reals:intervals}{intervalle} non dégénéré, $x_0$ son
milieu, $M$ arbitraire~: la \omterm{def:b1:topology:dense}{densité} fournit un point du graphe à
distance moindre que $\min\bigl(\frac{\abs I}{2}, 1\bigr)$ de $(x_0, M
+ 1)$, c'est-à-dire $x \in I$ avec $f(x) > M$~: non bornée sur $I$,
donc (question 8) discontinue en tout point et monotone sur aucun
\omterm{prop:b1:reals:intervals}{intervalle}~; et pour toute cible $y_0$, les points du graphe proches de
$(x_0, y_0)$ donnent des $f(x)$ arbitrairement proches de $y_0$ avec $x
\in I$~: $f(I)$ est \omterm{def:b1:topology:dense}{dense} dans $\R$. C'est la question 8 lue à
l'envers~: puisque tous les barreaux sont équivalents, une fonction
additive non linéaire doit les mettre en défaut \emph{tous}, partout.

\textbf{13.} Bien définie~: $m + n\sqrt2 = m' + n'\sqrt2$ impose
$(n - n')\sqrt2 = m' - m \in \Z$, donc $n = n'$ (sinon $\sqrt2 \in
\Q$) et $m = m'$. L'additivité sur $G$ est alors claire coordonnée par
coordonnée. Caractère non borné près de $0^+$~: fixons $\varepsilon \in
\intoo{0}{1}$ et $N \in \N$. Pour chaque $n$ fixé avec $\abs n
\leq N$, la condition $m + n\sqrt2 \in \intoo{0}{1}$ enferme $m$
dans un \omterm{prop:b1:reals:intervals}{intervalle} de longueur $1$~: au plus un entier $m$ par
$n$, donc au plus $2N + 1$ éléments de $G \cap \intoo{0}{1}$ vérifient
$\abs{f} \leq N$. Or $G \cap \intoo{0}{\varepsilon}$ est
infini ($G$ est \omterm{def:b1:topology:dense}{dense}, \cref{exo:b1:reals:9})~; il contient donc
un $g$ avec $\abs{f(g)} > N$~: $f$ n'est bornée sur aucun \omterm{def:b1:topology:open}{voisinage} à
droite de $0$. Prolonger $f$ en une fonction additive sur $\R$ revient
à choisir des valeurs de façon cohérente sur une famille de réels
$\Q$-linéairement indépendante et engendrant $\R$ sur $\Q$
--- une base de Hamel~; en produire une exige l'axiome du choix,
et aucune formule explicite ne peut le faire~: constructivement, nous
ne possédons le monstre que sur $G$.

\textbf{14.} $f(x) = f(\frac x2 + \frac x2) = f(\frac x2)^2 \geq
0$. Si $f(x_0) = 0$, alors $f(x) = f(x - x_0)f(x_0) = 0$ pour tout
$x$~: exclu. Donc $f > 0$ et $g = \ln \circ f$ est \omterm{def:b1:continuity:continuous}{continue}
(\cref{prop:b1:continuity:algebra}) avec $g(x + y) = g(x) +
g(y)$~: d'après l'\cref{exo:b1:continuity:11}, $g(x) = cx$, donc $f(x)
= \eu^{cx}$. Réciproquement chaque $\eu^{cx}$ convient~: les
\omterm{def:b1:structures:morphism}{morphismes} \omterm{def:b1:continuity:continuous}{continus} $(\R, +) \to (\R^*, \times)$ sont exactement les
exponentielles.

\textbf{15.} $g(u) = f(\eu^u)$ est \omterm{def:b1:continuity:continuous}{continue} et $g(u + v) =
f(\eu^u \eu^v) = g(u) + g(v)$~: $g(u) = cu$, et tout $x > 0$
s'écrit $x = \eu^u$ avec $u = \ln x$~: $f(x) = c\ln x$.

\textbf{16.} $h = \ln \circ f$ est \omterm{def:b1:continuity:continuous}{continue} sur
$\intoo{0}{+\infty}$ avec $h(xy) = h(x) + h(y)$~: d'après la question
15, $h(x) = c\ln x$, donc $f(x) = \eu^{c\ln x} = x^c$.

\textbf{17.} $x = y = 0$~: $f(0) = 2f(0) + f(0)^2$, donc $f(0)(1 +
f(0)) = 0$. Si $f(0) = -1$~: en posant $y = 0$, $f(x) = f(x) +
f(0) + f(x)f(0) = -1$ pour tout $x$~: la constante $f \equiv -1$
(qui vérifie effectivement l'équation). Sinon $f(0) = 0$~; $h
= 1 + f$ est \omterm{def:b1:continuity:continuous}{continue}, $h(0) = 1$, et
\[
h(x + y) = 1 + f(x) + f(y) + f(x)f(y) = h(x)\,h(y) :
\]
d'après la question 14, $h(x) = \eu^{cx}$, c'est-à-dire $f(x) =
\eu^{cx} - 1$ (le cas $c = 0$ donnant $f \equiv 0$). Liste complète~:
$f \equiv -1$ et $f(x) = \eu^{cx} - 1$, $c \in \R$.

\textbf{18.} $x = y = 0$~: $2f(0) = 4f(0)$, donc $f(0) = 0$. $x =
0$~: $f(y) + f(-y) = 2f(y)$, donc $f$ est paire. $y = x$~: $f(2x) =
4f(x)$. Récurrence en utilisant $(x, y) \to (nx, x)$~:
\[
f((n{+}1)x) = 2f(nx) + 2f(x) - f((n{-}1)x)
= (2n^2 + 2 - (n-1)^2) f(x) = (n+1)^2 f(x) .
\]
Puis $f(x) = f\bigl(q\cdot\frac xq\bigr) = q^2 f\bigl(\frac
xq\bigr)$ donne $f\bigl(\frac pq x\bigr) = \frac{p^2}{q^2}f(x)$~:
$f(r) = r^2 f(1)$ sur $\Q$ (la parité règle les signes). Les deux
fonctions \omterm{def:b1:continuity:continuous}{continues} $f$ et $x \mapsto f(1)x^2$ coïncident sur la
partie \omterm{def:b1:topology:dense}{dense} $\Q$, donc partout (question 24)~: $f(x) =
f(1)\,x^2$~; et tout $cx^2$ vérifie l'équation.

\textbf{19.} $g = f - f(0)$ est \omterm{def:b1:continuity:continuous}{continue}, $g(0) = 0$, et
vérifie l'équation de Jensen (les constantes se simplifient). En
prenant $y = 0$~: $g\bigl(\frac x2\bigr) = \frac{g(x)}{2}$. Alors pour
tous $x, y$~:
\[
\frac{g(x + y)}{2} = g\Bigl(\frac{x+y}{2}\Bigr)
= \frac{g(x) + g(y)}{2} ,
\]
donc $g$ est additive et \omterm{def:b1:continuity:continuous}{continue}~: $g(x) = cx$ (question 2), et
$f(x) = cx + d$ avec $d = f(0)$. Toutes les fonctions affines vérifient
Jensen~: la liste est complète.

\textbf{20.} Récurrence sur $n$. Pour $n = 0$~: $\lambda \in \{0,
1\}$, trivial. Supposons l'inégalité acquise pour tous les poids
$\frac{k}{2^n}$. Un poids $\lambda = \frac{k}{2^{n+1}}$ avec $k$
pair se ramène au niveau $n$~; pour $k = 2j + 1$, $\lambda$ est le
milieu de $\lambda_1 = \frac{j}{2^n}$ et $\lambda_2 =
\frac{j+1}{2^n}$. Avec $z_i = \lambda_i x + (1 - \lambda_i)y$~:
$\lambda x + (1-\lambda)y = \frac{z_1 + z_2}{2}$, donc
\[
f\bigl(\lambda x + (1{-}\lambda)y\bigr)
\leq \frac{f(z_1) + f(z_2)}{2}
\leq \frac{(\lambda_1 + \lambda_2)f(x) + (2 - \lambda_1 -
\lambda_2)f(y)}{2}
= \lambda f(x) + (1 - \lambda)f(y) .
\]

\textbf{21.} Fixons $x, y$. Les \omterm{def:b1:logic:map}{applications} $\lambda \mapsto f(\lambda
x + (1 - \lambda)y)$ et $\lambda \mapsto \lambda f(x) + (1 -
\lambda)f(y)$ sont \omterm{def:b1:continuity:continuous}{continues} sur $\intcc{0}{1}$ (composition et
algèbre, \cref{prop:b1:continuity:algebra}). L'inégalité
vaut sur les poids dyadiques, \omterm{def:b1:topology:dense}{denses} dans $\intcc{0}{1}$
(\cref{exo:b1:reals:8})~; pour $\lambda$ quelconque, prenons des
dyadiques $\lambda_n \to \lambda$ et passons à la limite
(\cref{thm:b1:continuity:seqchar} et \cref{thm:b1:seq:order})~:
l'inégalité de convexité vaut pour tout $\lambda \in
\intcc{0}{1}$ --- convexité au milieu plus \omterm{def:b1:continuity:continuous}{continuité} égale
convexité (la notion du \cref{ch:b1:derivative}).

\textbf{22.} Une fonction additive non linéaire $f$ vérifie
$f\bigl(\frac{x+y}{2}\bigr) = \frac{f(x) + f(y)}{2}$ exactement
(question 1 avec $r = \frac12$, puis additivité)~: elle est convexe au
milieu, et même affine au milieu. Si elle vérifiait l'inégalité de
convexité complète sur un \omterm{prop:b1:reals:intervals}{intervalle} $\intcc{x}{y}$, alors pour
$\lambda \in \intcc{0}{1}$~: $f(\lambda x + (1-\lambda)y) \leq
\max(f(x), f(y))$ --- majorée sur un \omterm{prop:b1:reals:intervals}{intervalle} non dégénéré,
donc linéaire d'après la question 9~: contradiction. La \omterm{def:b1:continuity:continuous}{continuité} de
la question 21 n'est donc pas un luxe~: sans elle, la convexité au
milieu ne contrôle que le squelette dyadique dénombrable, et le
\omterm{def:b1:continuity:continuous}{continu} qui s'étend entre ses points part en vrille.

\textbf{23.} $p(x) = \frac{x^2}{2}$ vérifie $p(x+y) = p(x) +
p(y) + xy$. Si $f$ est une solution \omterm{def:b1:continuity:continuous}{continue} quelconque, $g = f - p$
est \omterm{def:b1:continuity:continuous}{continue} et additive, donc $g(x) = cx$~:
\[
f(x) = \frac{x^2}{2} + cx , \qquad c \in \R ,
\]
et chacune de ces fonctions est solution~: la liste est complète.

\textbf{24.} Principe général~: si $u, v$ sont \omterm{def:b1:continuity:continuous}{continues} et
coïncident sur une partie \omterm{def:b1:topology:dense}{dense} $D \subseteq \R$, alors pour $x \in \R$
on choisit $d_n \in D$ avec $d_n \to x$
(\cref{prop:b1:topology:closureprops})~; $u(x) = \lim u(d_n) =
\lim v(d_n) = v(x)$. Soit maintenant $f$ \omterm{def:b1:continuity:continuous}{continue} et additive sur
le \omterm{def:b1:structures:subgroup}{sous-groupe} \omterm{def:b1:topology:dense}{dense} $G$. Fixons $x, y \in \R$ et prenons $g_n \to x$,
$g'_n \to y$ avec $g_n, g'_n \in G$~; alors $g_n + g'_n \to x + y$
et, par \omterm{def:b1:continuity:continuous}{continuité} séquentielle en $x + y$, en $x$ et en $y$~:
\[
f(x + y) = \lim f(g_n + g'_n) = \lim\bigl(f(g_n) +
f(g'_n)\bigr) = f(x) + f(y) :
\]
$f$ est additive sur $\R$ tout entier (donc linéaire, par la question
2).

\textbf{25.} (i) Pour $f$ additive~: linéaire $\iff$ \omterm{def:b1:continuity:continuous}{continue}
$\iff$ \omterm{def:b1:continuity:continuous}{continue} en un point $\iff$ monotone sur un \omterm{prop:b1:reals:intervals}{intervalle}
$\iff$ bornée (même d'un seul côté) sur un \omterm{prop:b1:reals:intervals}{intervalle}. (ii) La \omterm{def:b1:topology:dense}{densité}
du graphe dans le plan montre que l'échec n'est pas un défaut local
mais une explosion globale~: au-dessus de tout \omterm{prop:b1:reals:intervals}{sous-intervalle}, les
valeurs s'étalent sur $\R$ tout entier, de sorte que toute propriété de
régularité échoue partout à la fois. (iii) Le butin~: $cx$, $\eu^{cx}$,
$c\ln x$, $x^c$, $cx^2$, et $cx + d$ --- six caractérisations,
une seule méthode~: transporter l'équation vers celle de Cauchy,
démontrer le squelette sur $\Q$ par récurrence, remonter à $\R$ par
\omterm{def:b1:topology:dense}{densité} et \omterm{def:b1:continuity:continuous}{continuité}. (iv) La \omterm{def:b1:topology:dense}{densité} de $\Q$ (ou des dyadiques) a
porté les remontées de l'\cref{exo:b1:continuity:11} et de la question
21~; elle n'a rien porté à la question 13, car sans \omterm{def:b1:continuity:continuous}{continuité} les
valeurs ne se propagent pas d'une partie \omterm{def:b1:topology:dense}{dense} à son \omterm{def:b1:topology:closure}{adhérence} ---
la \omterm{def:b1:topology:dense}{densité} ne transfère l'information que le long de la \omterm{def:b1:continuity:continuous}{continuité}.
\end{solution}
